\begin{tikzpicture}[x=1cm,y=1cm,every node/.append style={anchor=center,inner sep=.3333em}]
\node[pFlow,text width=2.9cm] (a) at(0,2.4){\textbf{1. Stato corrente}\\$\theta^{(t)}$};
\node[pFlow,text width=3.4cm] (b) at(5.6,2.4){\textbf{2. Proposta}\\$\theta^*=\theta^{(t)}+\varepsilon$, $\varepsilon\sim N(0,s^2)$};
\node[pFlow,text width=3.4cm] (c) at(5.6,0){\textbf{3. Rapporto}\\$r=\dfrac{\pi(\theta^*\mid y)}{\pi(\theta^{(t)}\mid y)}$};
\node[pFlow,text width=3.2cm] (d) at(0,0){\textbf{4. Decisione}\\$\theta^{(t+1)}=\theta^*$ con prob. $\min(1,r)$};
\draw[pEdge] (a)--(b);\draw[pEdge] (b)--(c);\draw[pEdge] (c)--(d);\draw[pEdge] (d)--(a);
\node[pNote,text width=7.6cm] at(2.8,-1.65){La costante di normalizzazione si semplifica nel rapporto. Se la proposta è rifiutata, la catena ripete lo stato corrente (Metropolis e colleghi, 1953).};
\end{tikzpicture}
